\section{Resultados}

\subsection{FFT}

A Tabela~\ref{tab:fft} apresenta os tempos e speedups. A inclusão da
transferência reduziu o ganho, mas a GPU permaneceu mais rápida que a CPU.

\begin{table}[H]
\centering
\caption{Resultados da FFT para matriz $3000\times3000$.}
\label{tab:fft}
\begin{tabular}{lrr}
\toprule
Caso & Tempo (ms) & Speedup \\
\midrule
CPU & 137,556 & 1,000 \\
GPU & 1,969 & 69,875 \\
GPU + transferência & 9,015 & 15,259 \\
GPU residente & 2,953 & 46,582 \\
\bottomrule
\end{tabular}
\end{table}

\begin{figure}[H]
    \centering
    \includegraphics[width=\columnwidth]{fft_tempos.png}
    \caption{Tempos medidos para a FFT.}
    \label{fig:fft}
\end{figure}

As saídas foram numericamente equivalentes, com diferença absoluta máxima de
$1{,}85\times10^{-4}$.

\subsection{Mandelbrot}

A Tabela~\ref{tab:mandelbrot} mostra os resultados para grid $1000\times1000$
e 500 iterações.

\begin{table}[H]
\centering
\caption{Resultados do conjunto de Mandelbrot.}
\label{tab:mandelbrot}
\begin{tabular}{lrr}
\toprule
Implementação & Tempo (ms) & Speedup \\
\midrule
CPU serial & 210,842 & 1,000 \\
GPU vetorizada & 230,551 & 0,915 \\
Kernel CUDA & 4,677 & 45,079 \\
\bottomrule
\end{tabular}
\end{table}

\begin{figure}[H]
    \centering
    \includegraphics[width=\columnwidth]{mandelbrot_tempos.png}
    \caption{Comparação dos tempos do Mandelbrot.}
    \label{fig:mandelbrot-tempos}
\end{figure}

As matrizes produzidas pelas três implementações foram exatamente iguais:
nenhum pixel apresentou diferença.

\begin{figure}[H]
    \centering
    \includegraphics[width=0.92\columnwidth]{mandelbrot_gpu_kernel.png}
    \caption{Conjunto de Mandelbrot obtido pelo kernel CUDA.}
    \label{fig:mandelbrot}
\end{figure}
